\chapter{纯化}
\chapterkana{純化}
\epigraphbox{“纯粹而不杂，静一而不变。”}{——《庄子·刻意》}

\kaishou{操作表上有四列。前三列她认得。\par
最后一列她不认得，那一列是退出。\par
她按了前三列。}

\section{处置的手段}

上一章说到，新历的条文上行，月之都认出那是自己的底牌。认出之后要处置，而月之都手上能用的东西很少。

先说清楚它有什么。城里没有军队，这一点从工籍能看出来：工别四类，农作、清扫、制药、警备，四类都由月兔承担，没有一类是月之民自己做的。换句话说，这座城唯一能动手的力量，也是它们。\footnote{工别四类的全文见第三章。月兔之乱以前，警备一项从未用于月之民内部。}

城里的技术又是为无秽之身设计的。第二章说过，他们不病不老，章程里连医都没有一卷；一套没有医用条目的技术，落到会病会老的人身上，多半不成立。

当时唯一还通着的，是那部历法里的编码信道。信道本来有两个方向。上行是人按历作息，回传的是作息本身；下行是条文。月之都就是从下行这一头办的这件事。

第四十叶记这条规矩，写法很干：

\begin{archive}[内部参考]{档案 № ФВ-1971-22}
第四十叶。凡处置之令，皆由信道下行；一人手书；不署受令者。\\
页边另有一行，笔迹不同：「令出，不问其行。」
\end{archive}

手书的人是嫦娥。她从公元前二〇五〇年下地起，守着信道那一端，守了六十年。

条文怎么写，档案没有留底。能确定的只有一点：写法与其余律令一致，混在农时、祭祀、刑戮之间。后羿每日依历行事，读到哪一条就办哪一条。他不识字，没有理由把这一条与别条分开。

于是有一天，他对自己儿子伯封动了手。父子相搏，两个人都没有活下来。这件事的细节，传世文献只留下结果：伯封死了，后羿也死了。《天问》只问不答，问的是「何羿之射革，而交吞揆之」，没有一句交代那天发生了什么。

按本书的读法，后羿至死不知道自己在做什么，伯封至死认为被父亲背叛。这半句是笔者替他们说的，材料没有写；写下来的只有结果。

\section{纯狐}

这件事还有另一半，而另一半才是本章的题目。纯狐是后羿之妻，伯封之母，有穷氏的地上人。她不识字，不通术数，从未登上月球，也没有得到过月之民的任何技术。她手上有两样东西：一部旧历的抄本，和一处遗迹。

旧历是从夏那里流出来的，与后羿用过半辈子的那一部同源。遗迹在有穷氏居地之下，是月之民离开地球之前留下的设备，两万年间一直埋在原来的地方，没有人动过它。

传世文献里，纯狐最早出现在《天问》：「浞娶纯狐，眩妻爰谋。」\footnote{《楚辞·天问》。这一句把她记作寒浞之妻。本书取「后羿之妻」一说，取舍见本章末。} 那已经是战国人写的，去事一千七百年。她要报的仇，在后世的故事里被记成一个女人的妒忌；按本书的读法，那是一桩处置的后果落到了处置者家中。

\section{那台设备}

设备是本章唯一能摸到的东西，也是唯一有物证的。一九七二年，基金会对有穷氏故地做过一次地下测绘，编号 \texttt{ФВ-1972-06}，用电阻率成像，测出一个地下空腔，测线一共九条。

\begin{archive}[内部参考]{档案 № ФВ-1972-06}
地下空腔一处，深约十四米，长约六十米，宽约二十米。底部为金属反射面，当时完整，未见裂隙。
腔内壁有刻痕，分四列，字不可识，最右一列最密。
未取样，未进入。结论：不具备安全作业条件，建议维持封存。
\end{archive}

四列刻痕就是那张操作表。左起三列是定位、启动、监视，最右一列是退出顺序。设备停止之前，操作者须按最右一列逐步退出工作范围，人才出得来。

这台设备是月之民留在地面上的。留下来的不止这一处，规程残卷里另有一叶，列着当年置于地上的器若干所，只记方位与数，不记用途；有穷氏这一所在列。它两万年没有人维护，到前一九九〇年还能开动，只能说运气。

工作范围怎么划，规程里有一条通则，写得很短：「凡器，以执器者为心，周百步，皆在其内。」百步合一百三十八米，与空腔的尺度对得上。

这一条最要命的地方在于它不区分人。范围内的一切都处理，包括按表的人自己。

一九七二年测绘的时候，设备已经开不动了：金属反射面还在，内部没有可运转的结构，取样也只取到碎块。八十年代复查，空腔已经塌陷。今天能说的只有一句，它曾经能用，现在不能。

\section{三步}

纯狐怎么学会按那三列的，材料没有答案。能拼出来的只有一条路径，三样东西各管一段。

记号。她幼年在遗迹旁见过外壳上的刻痕，记住了形状。当时当地没有人能解释那是什么。

历法。夏所颁历法中另有一小段记号，与遗迹刻痕属于同一套。那一段正是上一章说的、删减时没有删掉的那一层：历不全的时候它算不出东西，看上去像几行旧文。

设备。遗迹底下那一台，提供执行的能力。

认得出记号，读得懂旧历里那一段，摸得到设备。三样凑齐，她就有了别人没有的东西。

此后她做了三步：划定范围，启动设备，没有退出。三步之中，前两步有物证，第三步只能从结果反推。

第三步是本章的关键，而它最没有材料。她没有退出，原因可能有三：不识字，读不懂最右那一列；来不及；根本不知道还有第四列。档案不写，本书也不替她选。能确定的只有结果：范围内的一切物质都被处理，她自己也在范围内。

人體被这样处理之后，不再是一个活的整体。她的死由此而来。

\section{之后}

被处理的那一片材料还在原处，位置没变，质量没变，只是不再属于嫦娥。这一句要慢慢说。蓬莱人的再生机制有一个前提：愈合的是「她自己的一部分」。前提一破，那一片便不再愈合，也不再腐烂，它已经不属于活体。

纯化没有停在最初那一片上。设备的残余作用没有被正确关闭，界线一直在外移。外移不快：起初只有一小部分身体受影响，其余仍是她自己的；界线之外的组织还能再生，每再生一次，新生的部分就以那一片为参照，随即被划进去。于是不能愈合、也不成原形的组织逐圈堆在边界上。

到接回的时候，形貌已经走了样。档案因此写作「酷似蟾蜍之物」。此后界线还在走，走到哪里为止，没有记录。

到撤离时为止，嫦娥仍然活着，住在隔离设施里，没有出来过。所谓隔离设施，是月面一处独立的封闭舱段，与六环城不连通。

神话把这整套事记成五个母题，与本书的读法对照如下：

\begin{table}[htbp]
\centering
\caption{五个母题与本书的读法}
\begin{tabular}{@{}ll@{}}
\toprule
神话 & 本书的读法\\
\midrule
奔月 & 接回 \\
窃药 & 西王母所赠之物被记成药，又被记成她窃取 \\
蟾蜍 & 界线外移之后的形貌 \\
捣药 & 月兔的本职，与嫦娥无关，神话记成替她治病 \\
囚禁 & 隔离 \\
\bottomrule
\end{tabular}
\tabsource{神话取自《淮南子·览冥训》及后世注疏。}
\end{table}

档案对这一整件事只有一行，末尾注「未解」。为什么写未解，可以推三层。他们认得出那是自家设备的用法，认不全为什么这样用。他们没有先例：那台设备两万年来只用来加工材料，加工活体这件事，在他们的认知里根本不存在。他们也不能写：承认一个地上人用他们的设备处理了月之民，等于承认核心装备外泄两次，第二次还用在人身上，这两件事都进不了档。于是档案里只剩一行，末尾加了两个字。

这一章有三处靠推：她怎样读懂那张表；她为什么没有退出；那一片为什么不合拢。第一处最要紧，而它一点材料也没有。三列刻痕、旧历里的一段、幼年记下的形状，三样东西怎么在一个不识字的人手里合成一次操作，本书拼不出完整的链。笔者能给的只是路径，路径不等于过程。

至于《天问》把她记作寒浞之妻，本书取后羿之妻一说，理由是伯封。按本书的读法，伯封必须是后羿之子，处置才落到他头上；若纯狐是寒浞之妻，这层关系就断了。两说都有文献，本书选的是能与档案对上的一条。选错的可能，笔者认。
